\section{Introduction}
\label{sec:intro}

Computer-use agents now operate real software the way people do: they read the
screen, move a mouse, and type. Recent benchmarks show these agents completing
tasks on real desktop applications \citep{xie2024osworld}, live websites
\citep{zhang2026clawbench}, command-line environments
\citep{merrill2026terminalbench}, and mobile devices
\citep{rawles2024androidworld}. As these agents move into products, they will be
asked to serve users in many languages. Whether they serve those users equally
well is an open measurement question.

Early evidence says they do not. A cross-lingual performance gap has now been
observed in every multilingual agent evaluation we are aware of.
\citet{yang2025macosworld} localize a macOS environment into five languages and
report an average degradation of 28.8\% in Arabic, the worst of their languages.
\citet{caciolai2026omnilingual} translate an agentic benchmark into ten
languages and find a universal gap of 8.8 to 18.4 pass@3 points that does not
close with model scale. \citet{wang2025xwebagentbench} report a lag of more than
20\% across fourteen languages in a simulated shopping environment, and
\citet{chen2025mprgui} find consistent non-English deficits in GUI perception
and reasoning. Across these studies, Arabic is consistently at or near the
bottom.

These findings motivate our work, but they also expose a measurement problem.
Every multilingual agent benchmark to date builds its non-English arms by
machine translation, and every one runs in a sandbox or simulation rather than
the real digital environment of the target users.
\citet{kim2026lilt} show that this is not a detail: auditing and correcting the
machine-translated arms of an agentic benchmark changed model scores by up to
32.7 points, meaning a substantial share of the reported multilingual gap was
measurement error rather than model failure. \citet{caciolai2026omnilingual}
reach a similar conclusion from the other direction, attributing 35\% of
observed failures to translation defects and 10\% to evaluation artifacts.
Prior work also says little about \emph{how} agents fail in Arabic. The
mechanisms reported so far are coarse: left-to-right click biases in mirrored
interfaces \citep{yang2025macosworld}, inflated token costs for non-Latin
scripts \citep{wang2025xwebagentbench,hofman2025maps}, and a behavioral shift
toward exploration over action \citep{caciolai2026omnilingual}. None of these
studies observes an agent inside a real Arabic-localized operating system, on
real Arabic-facing websites, or in any Arabic dialect.

No existing benchmark can therefore answer a question that is simple to state:
does the same agent complete the same real task as well for an Arabic-speaking
user as for an English-speaking one, and if not, what exactly goes wrong? The
answer requires holding the task constant, translating it with verified
fidelity, running both arms in matched real environments, and reading the
evidence behind every score. That is what we build.

We present \textbf{AraClawBench}, a paired English--Arabic benchmark for
computer-use agents. AraClawBench contains 221 tasks in three families: 171
tasks borrowed from established benchmarks \citep{xie2024osworld,
zhang2026clawbench, merrill2026terminalbench} and 50 native tasks drawn from the
digital life of an Arab user, from regional airlines and property portals to
charity donations and government-adjacent services. Every task carries a
human-translated Modern Standard Arabic instruction; a fixed subset carries four
additional dialect arms (Egyptian, Syrian, Qatari, and Sudanese). Both language
arms of a task run on matched virtual machines, English and Arabic-localized,
and the agent receives only screenshots and mouse and keyboard control. Each
task is graded by a deterministic checker that reads recorded evidence, never
the agent's self-report. Consequential actions --- payments, job applications,
public posts --- are performed against the real site but captured and blocked at
the network layer before they take effect, following \citet{zhang2026clawbench}.
Every English--Arabic score divergence is re-run and adjudicated from evidence
at $k{=}3$, and any cell whose environment, recording, or checker failed is
excluded as an evaluation error rather than scored as a model failure.

Using this benchmark we evaluate four frontier models: GPT-5.4, Claude Opus
4.8, Claude Sonnet 5, and Gemini 3.7 Flash.
\todo{Headline results paragraph once the audit lands: overall paired gap per
model with significance; whether the gap is universal across models; the
attribution decomposition (genuine failures vs.\ translation defects vs.\
evaluation artifacts); two or three leading mechanisms with counts; the
provenance split (human-translated vs.\ model-drafted arms); the dialect
result.}

Our contributions are:
\begin{enumerate}
    \item \textbf{The first paired English--Arabic computer-use benchmark.} 221
    real-environment tasks across desktop, web, and terminal, with
    human-translated Modern Standard Arabic arms, four dialect arms, and 50
    native Arab-world tasks. To our knowledge this includes the first agentic
    evaluation in any Arabic dialect.
    \item \textbf{A validity methodology for cross-language agent measurement.}
    Matched-environment paired cells; test-enforced frozen instructions with
    verbatim literal parity across arms; deterministic evidence-based checkers;
    capture-and-block grading of consequential actions; evaluation errors
    excluded rather than scored; $k{=}3$ adjudication of divergent pairs. This
    directly addresses the translation-induced measurement error documented by
    \citet{kim2026lilt}.
    \item \textbf{A controlled measurement of the Arabic gap.} A paired
    evaluation of four frontier models on the full corpus, with per-family
    significance testing, a gap attribution decomposition, and a
    mechanism-level failure taxonomy grounded in per-cell evidence.
    \item \textbf{Findings about the Arabic web itself.} Constructing the
    benchmark produced two ecosystem results: most of the Arabic web reachable
    by these tasks is machine-translated rather than natively Arabic, and
    phone-verification and CAPTCHA walls structurally exclude entire categories
    of regional services from agent evaluation. We report both with a full
    exclusion table.
\end{enumerate}
